
la série @@M0@@, où @@M1@@ est la classe de @@M2@@ dans @@M3@@. Son
noyau est le groupe @@M4@@ des séries formelles telles que~:
@@M5@@
on a donc un isomorphisme, fonctoriel en @@M6@@~:
@@M7@@

\textbf{Lemme 1.12.} \uline{Soit} @@M8@@ \uline{un foncteur sur} @@M9@@~; \uline{alors l'homomorphisme canonique}~:
@@M10@@
\uline{est un isomorphisme}.

L'homomorphisme @@M11@@ est défini comme étant le composé~:
@@M12@@
la dernière flèche étant définie par l'homomorphisme @@M13@@

On va construire un morphisme en sens inverse~:
@@M14@@
Pour cela, on va montrer que tout morphisme de foncteurs de @@M15@@ dans @@M16@@ se factorise par le morphisme canonique @@M17@@, ce qui équivaut à montrer qu'il s'annule sur @@M18@@ d'après 1.11.

Soient @@M19@@, @@M20@@~; on définit un homomorphisme de
